% Part: many-valued-logic
% Chapter: sequent-calculus
% Section: introduction

\documentclass[../../../include/open-logic-section]{subfiles}

\begin{document}

\olfileid{mvl}{seq}{int}

\olsection{مقدمة}

حساب التتابعيات في المنطق الكلاسيكي نسق !!{derivation} يجمع البساطة
إلى الفعالية. ولا يختص إمكان بنائه بذلك المنطق؛ فمتى عُرّف منطق متعدد
القيم بمصفوفة كانت مجموعة قيمها~$V$ منتهية، أمكن أن يُجعل له حساب
تتابعيات. ولبيان الطريق إلى ذلك، ننظر أولًا في دلالة التتابعية الكلاسيكية
وفي صورة قواعد الاستدلال عليها.

تذكر الآن أن التتابعية
\begin{align*}
    !A_1, \dots, !A_m & \Sequent !B_1, \dots, !B_n
\intertext{يمكن تفسيرها بالـ!!{formula}}
    (!A_1 \land \cdots \land !A_m) & \lif (!B_1 \lor \cdots \lor
!B_n)
\end{align*}
ومعنى ذلك أن !!^a{valuation}~$\pAssign{v}$ \emph{يحقق} التتابعية
$\Gamma \Sequent \Delta$ متى وفقط متى حصل أحد الأمرين:
$\pValue{v}(!A) = \False$ لبعض~$!A \in \Gamma$، أو
$\pValue{v}(!A) = \True$ لبعض~$!A \in \Delta$. وعلى هذا تكون
التتابعيات الابتدائية~$!A \Sequent !A$ محققة على كل حال؛ إذ لا يخلو
الأمر من~$\pValue v(!A) = \True$ أو~$\pValue v(!A) = \False$.

فيما يأتي قاعدتا الاستدلال الخاصتان بالشرط في~$\Log{LK}$، بعد حذف
الصيغ الجانبية~$\Gamma$ و$\Delta$:

\begin{defish}
    \Axiom$ \fCenter !A$
    \Axiom$ !B \fCenter $
    \RightLabel{\LeftR{\lif}}
    \BinaryInf$ !A \lif !B \fCenter $
    \DisplayProof
    \hfill
    \Axiom$ !A \fCenter !B$
    \RightLabel{\RightR{\lif}}
    \UnaryInf$ \fCenter !A \lif !B $
    \DisplayProof
\end{defish}

فلنقرأ القاعدتين على مقتضى هذه الدلالة. تفيد قاعدة
$\LeftR{\lif}$ أن اجتماع~$\pValue v(!A) = \True$ و
$\pValue v(!B) = \False$ يقتضي~$\pValue v(!A \lif !B) = \False$.
وتفيد قاعدة~$\RightR{\lif}$ أن حصول أحد الأمرين:
$\pValue v(!A) = \False$ أو~$\pValue v(!B) = \True$، يقتضي
$\pValue v(!A \lif !B) = \True$. بل يصح العكس في كل منهما،
فتكون الدلالة تكافؤًا. وليس هذا شأن قاعدتي~$\LeftR{\land}$ و$\RightR{\lor}$
في صورتيهما القياسيتين، غير أن لكل منهما صورة أخرى يتحقق فيها التكافؤ:

\begin{defish}
    \Axiom$!A, !B, \Gamma \fCenter \Delta$
    \RightLabel{\LeftR{\land}}
    \UnaryInf$!A \land !B, \Gamma \fCenter \Delta$
    \DisplayProof
    \hfill
    \Axiom$ \Gamma \fCenter \Delta, !A, !B$
    \RightLabel{\RightR{\lor}}
    \UnaryInf$ \Gamma \fCenter \Delta, !A \lor !B$
    \DisplayProof
\end{defish}

فإذا نقلنا هذا المعنى إلى منطق ذي~$n$ من القيم، كان حساب التتابعيات
ذا~$n$ من المواضع، لا موضعين فحسب: لكل قيمة صدق موضع يخصها.
ومثال ذلك المنطق الثلاثي القيم حيث~$V = \{\False, \Undef, \True\}$؛
فالتتابعية فيه على صورة~$\Gamma \mid \Pi \mid \Delta$. ويحققها
!!a{valuation}~$\pAssign v$ متى وفقط متى صدق أحد الشروط الآتية:
$\pValue{v}(!A) = \False$ لبعض~$!A \in \Gamma$، أو
$\pValue{v}(!A) = \True$ لبعض~$!A \in \Delta$، أو
$\pValue{v}(!A) = \Undef$ لبعض~$!A \in \Pi$.
ولذلك كانت التتابعيات الابتدائية~$!A \mid !A \mid !A$ محققة دائمًا.

\end{document}
